\begin{table}[!htb]
\centering
\caption{Síntese dos achados organizados pelas perguntas de pesquisa.}
\label{tab:revisao_sistematica_sintese}
\small
\begin{tabular}{>{\raggedright\arraybackslash}p{1.2cm}>{\raggedright\arraybackslash}p{5.9cm}>{\raggedright\arraybackslash}p{6.0cm}}
\toprule
RQ & Achado agregado & Implicação para a dissertação \\
\midrule
RQ1 & 48 estudos explicitam controle de qualidade ou artefatos; 11 tratam tecido/fundo como etapa operacional. & QC aparece como componente recorrente, mas com profundidade variável; por isso, deve ser tratado como contrato verificável do fluxo. \\
RQ2 & 29 estudos declaram separação por paciente; 33 mencionam normalização de cor. & A literatura reconhece particionamento e normalização como fontes de viés, mas nem sempre explicita a ordem operacional das etapas. \\
RQ3 & 56 estudos apresentam algum sinal de governança, repositório público, manifesto ou formato persistente. & A reprodutibilidade depende de rastrear lâmina, paciente, amostra e configuração, não apenas de publicar pesos de modelo. \\
RQ4 & 27 estudos discutem ensemble ou incerteza; 7 enfatizam pós-processamento. & Agregação de modelos é metodologicamente defensável quando a receita é definida em validação e congelada antes do teste. \\
\bottomrule
\end{tabular}
\normalsize
\end{table}
